\subsection{Feature importance (LRP)}\label{sec:lrpr}
\begin{table}[H]\centering\small
\caption{LRP sanity checks on the test set (15 models each). Corr/slope: Pearson correlation and slope of $\sum_j R_j$ vs.\ $f(x)$. Ratio: median $\sum_j R_j/f(x)$ for $|f(x)|>0.5$. LRP-0 vs.\ G$\times$I: maximum relative deviation $\max|R^{(0)}-G{\times}I|/\max|G{\times}I|$ over runs. Corr G$\times$I: correlation of the composite LRP maps with gradient$\times$input.}
\begin{tabular}{lcccccc}\toprule
Model (set) & Runs & Corr & Slope & Ratio & LRP-0 vs.\ G$\times$I & Corr G$\times$I\\\midrule
DFN (a) & 15 & 0.987 & 0.759 & 0.718 & 1.3e-04 & 0.984 \\
CNN1D (a) & 15 & 0.550 & 0.510 & 0.473 & 3.5e-02 & 0.166 \\
CNN1D (d) & 15 & 0.557 & 0.532 & 0.502 & 3.5e-04 & 0.144 \\
\bottomrule\end{tabular}\label{tab:lrpcons}\end{table}
LRP-0 matched gradient$\times$input to $10^{-4}$ for DFN and CNN1D (d), and to 3.5\% for CNN1D (a), where some near-zero denominators in the residual split amplify rounding. This confirms the implementation. Conservation is close for the DFN (corr 0.987, about 72--76\% of the output recovered). For CNN1D, about half of the output is recovered and the correlation is 0.55, because a large part of the relevance is absorbed by the folded BatchNorm biases. The composite $\gamma$ rule also changes the CNN maps substantially relative to gradient$\times$input (corr 0.14--0.17), whereas they are nearly identical for the DFN (0.984).

\begin{table}[H]\centering\scriptsize\setlength{\tabcolsep}{3pt}
\caption{LRP relevance per feature type (test rows, averaged over 15 models). abs: share of total $|R|$. signed: mean signed relevance over all test rows / over the top-100 test $k$-mers by affinity (standardised log-affinity units).}
\begin{tabular}{l|ccc|ccc|ccc}\toprule
& \multicolumn{3}{c|}{DFN (a)} & \multicolumn{3}{c|}{CNN1D (a)} & \multicolumn{3}{c}{CNN1D (d)}\\
Feature group & abs & signed all & signed top-100 & abs & signed all & signed top-100 & abs & signed all & signed top-100\\\midrule
one-hot & 0.347 & -0.002 & +0.331 & 0.489 & +0.072 & +0.525 & 1.000 & +0.146 & +1.063 \\
\texttt{avg} SHIFT & 0.013 & -0.000 & -0.007 & 0.022 & -0.000 & +0.010 & -- & -- & -- \\
\texttt{avg} SLIDE & 0.019 & +0.001 & +0.027 & 0.022 & +0.002 & +0.017 & -- & -- & -- \\
\texttt{avg} RISE & 0.019 & +0.020 & +0.121 & 0.020 & +0.011 & +0.053 & -- & -- & -- \\
\texttt{avg} TILT & 0.017 & -0.007 & -0.016 & 0.021 & +0.003 & +0.017 & -- & -- & -- \\
\texttt{avg} ROLL & 0.015 & -0.012 & -0.015 & 0.017 & -0.000 & +0.005 & -- & -- & -- \\
\texttt{avg} TWIST & 0.018 & -0.008 & +0.009 & 0.021 & +0.003 & +0.029 & -- & -- & -- \\
\texttt{fce} diagonal & 0.136 & -0.026 & +0.078 & 0.107 & +0.007 & +0.075 & -- & -- & -- \\
\texttt{fce} off-diagonal & 0.417 & +0.029 & +0.464 & 0.281 & +0.040 & +0.238 & -- & -- & -- \\
\bottomrule\end{tabular}\label{tab:lrpgroups}\end{table}

\begin{figure}[H]\centering\includegraphics[width=.85\textwidth]{lrp/groups_abs.png}
\caption{Share of total absolute LRP relevance per feature type, set (a).}\label{fig:lrpgroups}\end{figure}

\paragraph{Feature types.}
\begin{itemize}
\item \textbf{CNN1D (a):} one-hot carries 48.9\% of the absolute relevance, \texttt{fce} 38.9\% (off-diagonal 28.1\%, diagonal 10.8\%) and \texttt{avg} 12.2\%. Per channel, one-hot is the most used (0.031 per channel vs.\ 0.020 for \texttt{avg} and 0.019 for \texttt{fce}).
\item \textbf{DFN (a):} the DFN relies more on shape: \texttt{fce} carries 55.2\% (off-diagonal 41.7\%), one-hot 34.7\% and \texttt{avg} 10.1\%.
\item \textbf{\texttt{avg} parameters:} none dominates (1.3--2.2\% each). RISE has the largest positive signed contribution on high-affinity $k$-mers (+0.121 DFN, +0.053 CNN1D).
\item \textbf{\texttt{fce} terms:} the largest are TILT--ROLL (5.9\%), SHIFT--SHIFT (3.6\%) and RISE--RISE (3.3\%) for the DFN, and SLIDE--TILT (2.8\%), ROLL--TWIST (2.2\%) and SLIDE--ROLL (2.2\%) for CNN1D.
\item \textbf{Signed relevance} averages near zero over all rows. On the top-100 $k$-mers it is strongly positive for one-hot (+0.33 DFN, +0.52 CNN1D (a), +1.06 CNN1D (d)) and for \texttt{fce} off-diagonal (+0.46 DFN, +0.24 CNN1D): these features raise the prediction for strong binders.
\end{itemize}

\begin{figure}[H]\centering
\includegraphics[width=.49\textwidth]{lrp/pos_DFN_a.png}\includegraphics[width=.49\textwidth]{lrp/pos_CNN1D_a.png}
\includegraphics[width=.49\textwidth]{lrp/pos_CNN1D_d.png}
\caption{Mean $|R|$ of one-hot features per base position and of shape features per tetramer window.}\label{fig:lrppos}\end{figure}

\paragraph{Positions.}
\begin{itemize}
\item \textbf{CNN1D:} relevance peaks at the centre for one-hot (positions 5--6: 0.098/0.097 vs.\ 0.043/0.040 at the ends for set (a); 0.167/0.169 vs.\ 0.057/0.056 for (d)) and for shape (windows 3--5: 0.110--0.117 vs.\ 0.080 at windows 1 and 7). This follows from zero padding: edge positions enter fewer receptive fields.
\item \textbf{DFN:} the profile is U-shaped. One-hot relevance is highest at the end bases (0.21--0.22 vs.\ 0.10--0.11 at positions 4--5), and shape relevance is highest in windows 1 and 7 (0.44 vs.\ 0.35--0.39). The end bases are covered by only one tetramer window, so their identity is least determined by the shape block and the DFN reads it from one-hot.
\end{itemize}

\begin{figure}[H]\centering
\includegraphics[width=.9\textwidth]{lrp/heat_CNN1D_a.png}
\includegraphics[width=.9\textwidth]{lrp/heat_CNN1D_d.png}
\includegraphics[width=.9\textwidth]{lrp/heat_DFN_a.png}
\caption{Mean signed LRP relevance per nucleotide channel and base position: top-50 test $k$-mers by affinity (left) and all test rows (right).}\label{fig:lrpheat}\end{figure}

\begin{figure}[H]\centering
\includegraphics[width=.9\textwidth]{lrp/logo_CNN1D_d.png}
\includegraphics[width=.9\textwidth]{lrp/logo_CNN1D_a.png}
\includegraphics[width=.9\textwidth]{lrp/logo_DFN_a.png}
\caption{Relevance logos for the six highest-affinity test $k$-mers (letter height = signed relevance of the position).}\label{fig:lrplogo}\end{figure}

\paragraph{Motif check.} The highest-affinity $k$-mers do not contain a GATA site; they are ETS-like (ACTTCCGG / CCGGAAGT). Across all of \texttt{4.txt}:
\begin{itemize}
\item the 96 $k$-mers containing CCGGAAGT or ACTTCCGG have mean log affinity $-0.29$, vs.\ $-1.50$ for the rest;
\item $k$-mers containing GATA/TATC (5.3\%) are not enriched ($-1.48$ vs.\ $-1.50$);
\item the canonical WGATAR core AGATAA/TTATCT occurs in only 182 $k$-mers, all with count $\leq$9, so none reaches the count$\geq$10 training set.
\end{itemize}
In the test set, 2 of the top-100 $k$-mers contain GATA/TATC and 63 contain the ETS core GGAA/TTCC. In the logos, positive relevance falls on the flanking ACTT and GG of ACTTCCGG and on the GGAA run of CCGGAAGT, while the central CC gets little relevance. Measured as the share of a $k$-mer's one-hot $|R|$ that falls on its GGAA/TTCC core, the values for the top-100 $k$-mers are 0.342 (DFN), 0.426 (CNN1D (a)) and 0.447 (CNN1D (d)). The positional expectation for the same positions is 0.349, 0.427 and 0.433. Relevance therefore follows the motif in individual strong binders, but it is not concentrated on the 4-bp core beyond what the positional profile predicts.
